\section{Conclusion}
\label{sec:conclusion}

This paper presented the design, implementation, and empirical evaluation of an end-to-end retail demand forecasting system that trains nineteen forecasting models, spanning statistical, classical machine learning, deep learning, and intermittent-demand families, concurrently on an arbitrary uploaded time series, and selects among them using a robust, multi-criteria composite score rather than a single error metric. The system automatically detects relevant columns in heterogeneous uploaded files, applies a documented cleaning and feature-engineering pipeline, withholds deep learning models from datasets below an empirically justified minimum size rather than training them unreliably, and, when present, integrates real price and promotional exogenous features into the subset of models capable of using them. Beyond forecast accuracy, the system addresses the interpretability gap directly through native tree-based feature importance and a retrieval-augmented large-language-model explanation layer that is explicitly and architecturally constrained to remain consistent with the system's own deterministic metrics and labels, rather than reasoning independently over the underlying numbers.

Evaluated on a real retail transaction dataset aggregated to a 366-day daily series, twelve of the nineteen registered models completed successfully, with $k$-nearest-neighbor regression selected as the best model under the composite score despite extreme gradient boosting achieving a marginally lower root-mean-squared error, a divergence traced directly to the composite score's inclusion of stability and computation-time components alongside raw accuracy. The predominantly negative coefficients of determination observed across candidate models were reported transparently as a property of the specific, comparatively noisy evaluation dataset, consistent with the trustworthiness objective that motivated this work, rather than selectively omitted.

The central contribution of this paper is therefore architectural rather than algorithmic: it demonstrates that a broad, automatically adjudicated model ensemble, combined with explicit data-driven gating of unreliable model families and a grounded natural-language explanation layer, is a practical and extensible approach to making retail demand forecasting both more robust across heterogeneous datasets and more usable by non-specialist stakeholders. The limitations identified, fixed hyperparameters for the more complex model families, MAPE's known instability near zero-valued observations, restricted exogenous-feature scope, and the absence of hierarchical forecast reconciliation, together with the extensions outlined in Section~\ref{sec:future}, define a concrete roadmap for continued development of the system evaluated in this paper.
